\begin{table}[H]
  \centering
  \caption{问题三与问题二的衔接：两种信息边界下的全年结果（334 天）}
  \label{tab:q3-bridge}
  \small
  \begin{tabular}{llrrrr}
    \toprule
    方案 & 0:00 计划来源 & 全年费用/元 & 紧急购电量/kWh & 相对上一行节省/元 & 节省 95\% 区间/元 \\
    \midrule
    问题二方案（不调整） & 问题二日前计划 & 16,862,185.88 & 346,770.98 & -- & -- \\
    严格递进方案 & 沿用问题二合同 & 16,809,709.04 & 330,146.46 & 52,476.84 & $[-151{,}541.34,\;254{,}367.02]$ \\
    \textbf{联合方案（正式结果）} & 0:00 联合优化 & 15,773,614.72 & 465,083.60 & 1,036,094.32 & $[750{,}538.39,\;1{,}333{,}140.76]$ \\
    \bottomrule
  \end{tabular}
\end{table}

\begin{table}[H]
  \centering
  \caption{问题三消融实验（全年 334 天，锁定区间）：负荷预测口径、光伏插值与场景权重}
  \label{tab:q3-ablation}
  \small
  \begin{tabular}{lrrr}
    \toprule
    变体 & 全年费用/元 & 相对基准 & 紧急购电量/kWh \\
    \midrule
    基准（本文方案） & 15,773,614.72 & -- & 465,083.60 \\
    换成问题二正式七日均值负荷预测 & 18,950,991.51 & +20.14\% & 1,188,550.22 \\
    七日均值 + 光伏线性插值 & 19,361,699.29 & +22.75\% & 1,200,680.10 \\
    七日均值 + 插值 + 时间衰减场景 & 19,391,087.45 & +22.93\% & 1,215,461.50 \\
    本文预测 + 光伏线性插值 & 16,333,458.31 & +3.55\% & 471,377.91 \\
    本文预测 + 时间衰减场景 & 15,832,956.82 & +0.38\% & 480,888.65 \\
    本文预测 + 插值 + 时间衰减场景 & 16,400,321.42 & +3.97\% & 480,886.63 \\
    \bottomrule
  \end{tabular}
\end{table}

